% --------------------------------- Preamble -----------------------------------
\documentclass[11pt, letterpaper, twocolumn]{article}

\usepackage[T1]{fontenc} % use 8-bit fonts
\usepackage[margin=0.75in, columnsep=0.2in, % 3.4 in columns
    headsep=0.28in, footskip=0.42in]{geometry} % paper sizes
\usepackage{amsmath, amsfonts, amssymb} % improved math
\usepackage{bm} % bold math (must be after amsmath)
\usepackage{upquote} % upright single quotes in verbatim
\usepackage{pgfplots} % includes tikz, xcolor, graphicx
\usepackage[colorlinks=true, allcolors=gray]{hyperref}

\raggedbottom % do not force stretching text to fill the page
\setlength{\emergencystretch}{\hsize} % avoid overfull lines
\overfullrule=1mm % black band to show overfull lines
\pdfsuppresswarningpagegroup=1 % remove pdf image warnings
\frenchspacing % consistent spacing between words
\newcommand{\ul}[1] % vectors
    {{}\mkern1mu\underline{\mkern-1mu#1\mkern-1mu}\mkern1mu}
\DeclareSymbolFont{euler}{U}{eur}{m}{n} % symbol font
\DeclareMathSymbol{\PI}{\mathord}{euler}{25} % up pi
\pgfplotsset{compat=newest} % use newest pgfplot version
\newcommand{\EE}[2]{\ensuremath{{{#1}\!\times\!10^{#2}}}}
% ------------------------------------------------------------------------------

\title{Kalman Filter Tuning and Observability Analysis}
\author{}
\date{}

\begin{document}
\maketitle

\section*{Description}

This project extends a linear Kalman Filter implementation by evaluating estimator
consistency and system observability. In addition to state estimation, three
performance metrics are studied: the mean squared error (MSE), the average normalized
error estimate squared (ANEES), and the average normalized innovation squared (ANIS).
The project also examines how observability changes under modifications to the 
measurement and dynamics models.

\subsection*{System Model}

The system is the same linear position-velocity model used in the previous Kalman
Filter project. The state is

\[
x = \begin{bmatrix}
    p_x \\
    p_y \\
    v_x \\
    v_y
\end{bmatrix}
\]
with linear dynamics and linear measurements. Process and measurement noise are 
modeled as zero-mean Gaussian random variables with covariances $Q$ and $R$.

Truth data are generated using fixed baseline values of $Q$ and $R$, while the
estimator is run using scaled versions of these matrices to study the effect
of covariance mismatch.

\section*{Consistency Metrics}

Three metrics are used to evaluate estimator performance.

\textbf{Mean Squared Error (MSE):}
\[
\mathrm{MSE} = \frac{1}{K} \sum_{k=1}^{K} \ul{e}_k^{\top}\ul{e}_k
\]
where
\[
\ul{e}_k = \ul{x}_k - \hat{\ul{x}}_k
\]
is the state estimation error.

\textbf{Average Normalized Error Estimate Squared (ANEES):}
\[
\mathrm{ANEES} = \frac{1}{KN} \sum_{k=1}^{K} \ul{e}_k^{\top}\hat{\mathbf{P}}^{-1}\ul{e}_k
\]
This metric computes the actual estimation error to the filter's predicted
covariance. For a consistent estimator, the ANEES should approach 1.

\textbf{Average Normalized Innovation Squared (ANIS):}
\[
\mathrm{ANIS} = \frac{1}{KM} \sum_{k=1}^{K} \ul{r}^{\top}\hat{\mathbf{S}}^{-1}\ul{r}
\]
where
\[
\ul{r} = \ul{z} - \mathbf{H}\hat{\ul{x}}
\]
is the innovation and
\[
\mathbf{S} = \mathbf{H}\hat{\mathbf{P}}\mathbf{H}^{\top} + \mathbf{R}
\]
is its covariance. Like the ANEES, the ANIS should also approach 1 for a 
well-tuned estimator.

\section*{Covariance Scaling Study}

The truth data and measurements are generated using fixed baseline covariance
matrices $Q$ and $R$. The estimator is then run over 49 combinations of scaled
covariance matrices using the factors
\[
\left\{ \frac{1}{16}, \frac{1}{9}, \frac{1}{4}, 1, 4, 9, 16 \right\}
\]

The process noise covariance scaling increases down the rows of the analysis
matrices, and the measurement noise covariance scaling increases across the 
columns. For each combination, the MSE, ANEES, and ANIS metrics are computed
and stored in $7 \times 7$ matrices.

This study shows how sensitive the estimator is to incorrect assumptions 
about the process and measurement uncertainty.

\section*{Observability Analysis}

The observability of the system is analzyed using the observability matrix
\[
O = \begin{bmatrix}
    \mathbf{H} \\
    \mathbf{HF} \\
    \vdots \\
    \mathbf{H}\mathbf{F}^{l-1}
\end{bmatrix} \quad
l = \left\lceil \frac{N}{M} \right\rceil
\]

The singular value decomposition (SVD) of the observability matrix is then computed:
\[
O = USV^{\top}
\]

The singular values indicate the strength of observability in different directions of
the state space, while the rows of $V^{\top}$ give the corresponding observability
vectors. Zero or near-zero singular values indicate non-observable or weakly 
observable directions.

\section*{Modified Systems}

Two system modifications are studied.

\subsection*{Measurement Modification}

The measurement model is changed so that one position-related measurement is replaced
by a velocity measurement component. The observability matrix, singular values, and
observability vectors are then recomputed. The modified H matrix is
\[
\mathbf{H} = \begin{bmatrix}
    \cos(\theta) & \sin(\theta) & 0 & 0 \\
    0 & 0 & 1 & 0
\end{bmatrix}
\]

\subsection*{Dynamics Modification}

Keeping the modified measurement model, the dynamics are further changed so that the
state coordinates rotate at a constant rate. This alters the coupling between the
states and changes the observability of the system. The resulting observability
matrix and singular values demonstrate that both the dynamics model and the 
measurement model joint determine which state directions are observable.
The modified $\Phi$ matrix is
\[
\Phi = \begin{bmatrix}
    0 & \gamma & 1 & 0 \\
    -\gamma & 0 & 0 & 1 \\
    0 & 0 & 0 & \gamma \\
    0 & 0 & -\gamma & 0
\end{bmatrix}
\]
where $\gamma = \frac{4\pi}{180}$ corresponds to a rotation of the coordinate 
sytem of the state space at the rate of 4 degrees per second clockwise.

\section*{Results}

Figures~\ref{fig:mse}, \ref{fig:anees}, and \ref{fig:anis} show the MSE, ANEES,
and ANIS values across all 49 covariance-scaling combinations.

\begin{figure}[h]
    \centering
    \includegraphics[width=\linewidth]{../figures/fig_mse.pdf}
    \caption{MSE under process and measurement covariance scaling}
    \label{fig:mse}
\end{figure}

\begin{figure}[h]
    \centering
    \includegraphics[width=\linewidth]{../figures/fig_anees.pdf}
    \caption{ANEES under process and measurement covariance scaling}
    \label{fig:anees}
\end{figure}

\begin{figure}[h]
    \centering
    \includegraphics[width=\linewidth]{../figures/fig_anis.pdf}
    \caption{ANIS under process and measurement covariance scaling}
    \label{fig:anis}
\end{figure}

\pagebreak

Figures~\ref{fig:unmodified}--\ref{fig:modified} show the observability
matrices and singular values for the modified and unmodified systems.

\begin{figure}[h]
    \centering
    \[ \begin{aligned}
        s &= \{ 1, 1, 1, 1 \} \\
        \bm{V}^\top &= \begin{bmatrix}
            -1 & 0 & 0 & 0 \\
            0 & 1 & 0 & 0 \\
            0 & 0 & -1 & 0 \\
            0 & 0 & 0 & 1
        \end{bmatrix}
    \end{aligned} 
    \]
    \caption{Observability matrix and singular values for unmodified system}
    \label{fig:unmodified}
\end{figure}

\begin{figure}[h]
    \centering
    \[ \begin{aligned}
        s &= \{ 1.31, 1, 0.54, 0 \} \\
        \bm{V}^\top &= \begin{bmatrix}
            0 & 0 & 0.92 & 0.38 \\
            0.71 & 0.71 & 0 & 0 \\
            0 & 0 & -0.38 & 0.92 \\
            -0.71 & 0.71 & 0 & 0
        \end{bmatrix}
    \end{aligned} 
    \]
    \caption{Observability matrix and singular values for measurement-only modification}
    \label{fig:measonly}
\end{figure}

\begin{figure}[h]
    \centering
    \[ \begin{aligned}
        s &= \{ 1.31, 1, 0.55, 0.01 \} \\
        \bm{V}^\top &= \begin{bmatrix}
            0.03 & -0.03 & -0.92 & -0.38 \\
            0.71 & 0.71 & 0 & 0 \\
            0.06 & -0.06 & 0.39 & -0.92 \\
            0.7 & -0.7 & 0 & 0.1
        \end{bmatrix}
    \end{aligned} 
    \]
    \caption{Observability matrix and singular values for measurement and dynamics modification}
    \label{fig:modified}
\end{figure}

\pagebreak

\section*{Observations}

\subsection*{MSE}

The highest MSE values were the ones with Q and R having opposing end
scale factors. For example Q scaled by 1/16 and R scaled by 16 leads to
a large MSE because our filter will almost solely follow the dynamics.
Similarly, Q scaled by 16 and R scaled by 1/16 will have the opposite
effect on the filter, essentially following only the measurements. Both
of these cases and cases where Q and R are scaled similarly apart produce
large MSE's because our dynamics and measurements are noisy by themselves.
It then follows logically that when Q and R are scaled the same amount, it
leads to the lowest (and best) MSEs. 

One interesting thing of note is that
our MSEs tend to be slightly better, in this scenario, when Q and R are
scaled the highest (but still the same amount). This is actually a minor
bug in our setup, which is explained in two ways. The first is that the 
covariance matrix P evolves according to the discrete-time Riccati recursion.
The solution to this equation requires an infinite time horizon. So, if we 
increased the number of steps in our simulation, our MSE would be equal along
the big diagonal. The second explanation is that when we scale Q by
some value k, we should also scale the prior P0 by the same amount. Scaling
P0 by the same amount as Q fixed this behavior along the diagonal, as expected.

\subsection*{ANEES}

The highest ANEES values are in the top left corner, where both Q and R
are scaled to be very low. This leads to a very small covariance matrix
which harshly penalizes our larger errors. A high ANEES score means our
filter is overconfident, which makes sense with Q and R values lower than
truth. Similarly, the lowest ANEES values are in the lower right corner,
where both Q and R are scaled to be very high. This leads to a large covariance
matrix which forgives even the largest errors. A low ANEES score means our 
filter is underconfident, which makes sense with Q and R values higher
than truth. Lastly, our closest values to 1 are in the middle, where the
filter Q and R most accurately represent the truth Q and R.

\subsection*{ANIS}

Similar to ANEES, the ANIS values are highest in the top left corner,
and the lowest in the bottom right corner. The main difference with ANIS
is that the score is more sensitive to changes in R. This is because 
ANIS directly deals with R when it calculates difference from expected 
measurement. The fact that the R noise is higher than the Q noise in 
generation and in the baseline estimation is also a factor. We again
notice that our values closest to 1 are at the positions closest to 
the truth Q and R, with R being more sensitive here too.

\subsection*{Observability}

For both the unmodified case and the case where the measurements and
dynamics were modified, we notice that all states in the system are
observable, though the unmodified case had ``more'' observable states.

In the case where only the measurements are modified, we note one 
linear combination of states that is ``unobservable''. The unobservable
direction is $-p_x$ + $p_y$ as shown by $s_4 = 0$ and row 4 of 
$V^T$ = $-sin\theta$, $cos\theta$, 0, 0.

\section*{Discussion}

The covariance scaling study shows that estimator performance depends
strongly on the assumed values of $Q$ and $R$. When these covariances
are poorly tuned, the filter may still produce estimates, but the resulting 
uncertainty no longer matches the true estimation error. This is reflected
by ANEES and ANIS values that deviate from 1.

The observability analysis shows that accurate estimation is not determined
by the measurement model alone. The dynamics also influence which directions
in state space can be reconstructed from the available measurements. 
Modifying the measurement model can reduce observability, while modifying
the dynamics can restore coupling that improves it.

Together, these results highlight two central ideas in state estimation:
a filter must be both well tuned and applied to a system whose states
are sufficiently observable.

\end{document}